\section{A bounded seam observable on one spatial cube}\label{cl:sec:bounded-seam}

Consider the two temporal cuts in Section~\ref{cl:sec:replica}. At each cut
choose a matching sign $\sigma,\tau\in\{-1,1\}$, with equal independent
prior probabilities, as in the Anaxia replica-seam construction.%
\footnote{The construction workspace is linked in Data and Code Availability.}
Conditional on the signs, use the corresponding
Wilson weight on both sheets and integrate all links against Haar measure.
The probability below is taken after jointly normalising the four
weighted ensembles. The argument applies to each Wilson group in
Proposition~\ref{sc:prop:wilson}: relabelling the sheets preserves the
plaquette traces and inactive Haar bundles have integral one.

For a fixed spatial box $s$, write $A=Z(s,T)^2$ and $B=Z(s,2T)$.
Each equal-sign corner has weight proportional to $A$, and each
opposite-sign corner has weight proportional to $B$. Thus the probability
of different matchings and the correlation of the signs are
\begin{equation}\label{cl:eq:seam-disagreement}
 p_{\rm diff}=\frac{B}{A+B},\qquad
 \mathbb E[\sigma\tau]=1-2p_{\rm diff}=\tanh\bigl(\psi_s(T)/2\bigr).
\end{equation}
The signs have mean zero. This identity concerns the Wilson-weighted
probability. The prior probability of disagreement is $1/2$.

\begin{theorem}[A single seam-disagreement test]\label{cl:thm:seam-disagreement}
Fix an integer $T\ge3$ and one bare coupling. If
\begin{equation}\label{cl:eq:two-fifths}
 p_{\rm diff}\ge 2/5 \quad\hbox{on the spatial cube } (24T)^3,
\end{equation}
then, writing $q=3^{-1/7}$, every spatial cube of integer side $n\ge24T$
satisfies
\[
 r_1(n^3)\le q^{1/T},\qquad
 \Dexc{n^3}{n}\le q^{n/T}.
\]
Suppose there is $t_0>0$ such that for all sufficiently large $p$ and,
for each such $p$, all sufficiently large $N$ in its mesh sequence,
\eqref{cl:eq:two-fifths} holds at the actual coupling $\beta_{N,p}$ for
some integer $3\le T\le Nt_0$. The physical gap in a compatible Wightman
limit under hypothesis \textup{(W)} is then at least $\log(3)/(7t_0)$.
\end{theorem}
\begin{proof}
Put $x_j=r_j(s)^T$ for the excited eigenvalues, counted with multiplicity,
$D=\sum_{j\ge1}x_j$ and $E=\sum_{j\ge1}x_j^2$. Equation
\eqref{cl:eq:two-fifths} yields $(1+D)^2/(1+E)=A/B\le3/2$.
Since $E\le D$, this first gives $D\le1/2$. Since also $E\le D^2$,
it gives $D^2-4D+1\ge0$. Hence $D\le2-\sqrt3<2/7=:d$.

Let $T_i=8^iT$, and put $D_i=\Dexc{(24T_i)^3}{T_i}$.
Raise the transfer operator to a higher power on the \emph{same} box
before enlarging that box. Positivity of the excited weights gives
$\Dexc{(24T_i)^3}{8T_i}\le D_i^8$. At time $8T_i$ apply
Theorem~\ref{sc:thm:cube-control} with initial spatial sides $24T_i$.
For each integer $24T_i\le n\le192T_i$ the enlargement factor is at
most $(192/16)^3=1728$. It follows that
\begin{equation}\label{cl:eq:eightfold-bound}
 \Dexc{n^3}{8T_i}\le(1+D_i^8)^{1728}-1.
\end{equation}
The endpoint $n=192T_i$ gives $D_{i+1}$ on the left.
For $0\le x\le d$ one has $1728d^8<1/8$ and therefore
\[
 (1+x^8)^{1728}-1
 \le e^{1728x^8}-1
 \le \frac{1728x^8}{1-1728d^8}\le2048x^8 .
\]
Also $2048d^7=262144/823543<1/3$. Induction gives
$D_i\le d q^{8^i-1}$: if this holds at $i$, the right-hand side of
\eqref{cl:eq:eightfold-bound} is at most
$d(2048d^7)q^{8(8^i-1)}\le d q^{8^{i+1}-1}$.

For any $n\ge24T$ take $i$ with $24T_i\le n<192T_i$. Since $d<q$,
\eqref{cl:eq:eightfold-bound} gives
$r_1(n^3)^{8T_i}\le\Dexc{n^3}{8T_i}\le q^{8^{i+1}}$.
Taking the $8T_i$-th root gives the claimed eigenvalue bound. Multiplying
each term of the excited trace by the remaining power of its eigenvalue
gives $\Dexc{n^3}{n}\le q^{8^{i+1}}q^{(n-8T_i)/T}=q^{n/T}$.
On $n=NL_p$, once $L_p\ge24t_0$, these bounds give
$r_1\le e^{-\Delta/N}$ and $\Dexc{n^3}{n}\le e^{-\Delta L_p}$ with
$\Delta=\log(3)/(7t_0)$. The first gives the exponential bound on
centered transfer matrix elements and the second removes the excited
thermal trace in the large-volume limit. The thermal comparison and
spectral argument in the proof of Theorem~\ref{sc:thm:continuum}
now apply under the same hypothesis \textup{(W)}.
\end{proof}

An elementary bound for the actual Wilson ensemble illustrates the scale
of \eqref{cl:eq:two-fifths}. Write $b=N_c^2\beta$, so each plaquette
term in the exponent lies in $[-b,b]$. Switching one cut changes six
terms per spatial site (three directions on each sheet). The exponent
changes by at most $12bV$ on the common link space. Integrating gives
$A/B\le e^{12bV}$ and hence $p_{\rm diff}\ge(1+e^{12bV})^{-1}$.
This lower bound tends to zero on expanding weak-coupling cubes.
Improving it to $2/5$ requires control of the integrated response
between the cuts.

This test retains information about all excited states. For illustration,
take the even occupation numbers of one free massless lattice oscillator
of momentum $2\pi/(24T)$, with the vacuum energy subtracted. Its energy
spacing before the even projection is
$E_1=2\operatorname{arsinh}(\sin(\pi/(24T)))\le\pi/(12T)$.
Writing $z=e^{-2TE_1}$, the two traces are $(1-z)^{-1}$ and
$(1-z^2)^{-1}$. Formula~\eqref{cl:eq:seam-disagreement} gives
$1-2p_{\rm diff}=z\ge e^{-\pi/6}>1/5$. Even this single oscillator
fails \eqref{cl:eq:two-fifths}. The interacting Wilson probability along
the tuning requires an independent estimate. No estimate of that
probability follows from the prior on the matching signs or from the
one-link vacuum comparison.
